% PDF page 346
\source{gilbarg-trudinger:page-0346}

and hence by (ii)

\[
\frac{w^-(x)-w^-(x_0)}{|x-x_0|}\leq \frac{u(x)-u(x_0)}{|x-x_0|}\leq \frac{w^+(x)-w^+(x_0)}{|x-x_0|}.
\]

Consequently, provided they exist at $x_0$, the normal derivatives of $w^\pm$ and $u$ satisfy

\begin{equation}
\frac{\partial w^-}{\partial v}(x_0)\leq \frac{\partial u}{\partial v}(x_0)\leq \frac{\partial w^+}{\partial v}(x_0).
\tag{14.3}
\end{equation}

We call the functions $w^\pm$ respectively \emph{upper} and \emph{lower} barriers at $x_0$ for the operator
$Q$ and function $u$. Their existence at all points $x_0\in\partial\Omega$, with uniformly bounded
gradients, implies the desired boundary gradient estimate for $u$ satisfying $Qu=0$
in $\Omega$.

    Before setting about the construction of barriers, it is convenient to have certain
transformation formulae at our disposal. First, let $I$ be some interval in $\mathbb{R}$ and set
$u=\psi(v)$ where $\psi\in C^2(I)$, $\psi'\neq 0$ on $I$. We then have, for $v(x)\in I$,

\begin{equation}
Qv=\psi' a^{ij}D_{ij}v+\frac{\psi''}{(\psi')^2}\mathcal E+b
\tag{14.4}
\end{equation}

where the arguments of $a^{ij}$, $\mathcal E$ and $b$ are $x$, $u=\psi(v)$ and $Du=\psi' Dv$. Next, let us
set $u=v+\varphi$ for some $\varphi\in C^2(\overline{\Omega})$. Then we have for $x\in \Omega$,

\begin{equation}
\widetilde{Q}v=Qu=a^{ij}D_{ij}v+a^{ij}D_{ij}\varphi+b,
\tag{14.5}
\end{equation}

where the arguments of $a^{ij}$ and $b$ are $x$, $u=v+\varphi$ and $Du=Dv+D\varphi$. Defining
the function $\mathcal F$ by

\begin{equation}
\mathcal F(x,z,p,q)=a^{ij}(x,z,p)(p_i-q_i)(p_j-q_j),
\tag{14.6}
\end{equation}
\[
(x,z,p,q)\in \Omega\times \mathbb{R}\times \mathbb{R}^n\times \mathbb{R}^n,
\]

we see that for the transformed operator $\widetilde{Q}$ in (14.5) we have

\begin{equation}
\widetilde{Q}(x,v,Dv)=\mathcal F(x,u,Du,D\varphi).
\tag{14.7}
\end{equation}

Although the formula (14.5) has been implicitly used in Chapter 13 to effect a
subtraction of boundary values, the explicit relation (14.7) is important for our
purposes here. As will be evident from the considerations of this chapter, the
formulae (14.4) and (14.5) foreshadow to some extent the nature of the structure
conditions required for our barrier constructions. We shall assume throughout
this chapter that the operator $Q$ is elliptic in the domain $\Omega$.


% PDF page 347
\source{gilbarg-trudinger:page-0347}


14.1. General Domains

14.1. General Domains

We commence with a barrier construction that is applicable to arbitrary smooth
domains. Suppose that $\Omega$ satisfies an exterior sphere condition at a point
$x_0 \in \partial\Omega$ so that there exists a ball $B=B_R(y)$ with
$x_0 \in \overline{B} \cap \overline{\Omega} = \overline{B} \cap \partial\Omega$.
Let us define the distance function $d(x)=\operatorname{dist}(x,\partial B)$ and
set $w=\psi(d)$ where $\psi \in C^2[0,\infty)$ and $\psi'>0$. By virtue of
formula (14.4), we have for any $u \in C^1(\overline{\Omega}) \cap C^2(\Omega)$

\begin{equation}\tag{14.8}
\overline{Q}w
=a^{ij}(x,u(x),Dw)D_{ij}w+b(x,u(x),Dw)
\end{equation}
\[
=\psi' a^{ij}D_{ij}d+b+\frac{\psi''}{(\psi')^2}\,\mathcal{E}
\]
\[
\leq \frac{n-1}{R}\,\psi'\,\Lambda+b+\frac{\psi''}{(\psi')^2}\,\mathcal{E};
\]

the last inequality follows since

\[
D_{ij}d=|x-y|^{-3}\bigl(|x-y|^2\delta_{ij}-(x_i-y_i)(x_j-y_j)\bigr).
\]

We now suppose that $Q$ satisfies a structure condition. Namely, we assume the
existence of a non-decreasing function $\mu$ such that

\begin{equation}\tag{14.9}
|p|\Lambda+|b|\leq \mu(|z|)\mathcal{E}
\end{equation}

for all $(x,z,p)\in \Omega\times \mathbb{R}\times \mathbb{R}^n$ with
$|p|\geq \mu(|z|)$. Using the condition (14.9) in (14.8), we obtain

\begin{equation}\tag{14.10}
\overline{Q}w \leq \left(\frac{\psi''}{(\psi')^2}+\nu\right)\mathcal{E}
\end{equation}

provided $\psi' \geq \mu=\mu(M)$, where $\nu=(1+(n-1)/R)\mu$, $M=\sup_\Omega |u|$.
Consider now the function $\psi$ given by

\begin{equation}\tag{14.11}
\psi(d)=\frac{1}{\nu}\log(1+kd), \qquad k>0,
\end{equation}

and the neighborhood $\mathcal{N}=\mathcal{N}_{x_0}=\{x\in\overline{\Omega}\mid
d(x)<a\}$, $a>0$. Clearly $\psi''=-\nu(\psi')^2$ in $\mathcal{N}$. Furthermore,

\begin{equation}\tag{14.12}
\psi(a)=\frac{1}{\nu}\log(1+ka)=M \qquad \text{if } ka=e^{\nu M}-1
\end{equation}


% PDF page 348
\source{gilbarg-trudinger:page-0348}

and

\begin{equation}
\psi'(d)=\frac{k}{v(1+kd)}\geq \frac{k}{v(1+ka)} \quad \text{in } \mathcal{N}\cap \Omega
\tag{14.13}
\end{equation}
\[
=\frac{k}{v\,e^{vM}}
\]
\[
\geq \mu \qquad \text{if } k\geq \mu v\, e^{vM}.
\]

Consequently, if $k$ and $a$ are chosen to satisfy the relations

\begin{equation}
k\geq \mu v\, e^{vM}, \qquad ka=e^{vM}-1,
\tag{14.14}
\end{equation}

the function $w^+=\psi(d)$ is an upper barrier at $x_0$ for the operator $\overline{Q}$ and the function
$u$ provided $u=0$ on $\mathcal{N}\cap \partial\Omega$. Similarly the function $w^-=-\psi(d)$ is a corresponding
lower barrier. Hence if also $Qu=0$ in $\Omega$, we obtain from (14.3) the estimate

\begin{equation}
|Du(x_0)|\leq \psi'(0)=\mu e^{vM} \qquad \text{if equality holds in (14.14).}
\tag{14.15}
\end{equation}

    Since the estimates to follow throughout this chapter are derived by essentially
the same argument as above but with the surface $\partial B$ replaced by other surfaces, it
is worth illustrating the situation by the accompanying Figure 2.

\begin{center}
Figure 2
\end{center}

\noindent\textbf{We now extend the estimate (14.15) to non-zero boundary values. Let $\varphi\in C^2(\overline{\Omega})$}
and suppose that $u=\varphi$ on $\partial\Omega$. Then we require the transformed operator given by
(14.5) to satisfy the structure condition (14.9). It therefore suffices that

\begin{equation}
\left(|p-D\varphi|+|D^2\varphi|\right)\Lambda+|b|\leq \bar{\mu}(|z|)\mathcal{F}(x,z,p,D\varphi)
\tag{14.16}
\end{equation}


% PDF page 349
\source{gilbarg-trudinger:page-0349}




for all $(x, z, p)\in \Omega\times\mathbb{R}\times\mathbb{R}^n$ such that $|p-D\varphi|\geq \mubar(|z|)$ for some non-decreasing function $\mubar$. Since

\begin{equation}\tag{14.17}
\begin{aligned}
\F(x, z, p, q)&=a^{ij}(x, z, p)(p_i-q_i)(p_j-q_j)\\
&\geq \frac12 \mathcal{E}^0-a^{ij}q_iq_j \qquad \text{by the Schwarz inequality,}\\
&\geq \frac12 \mathcal{E}^0-A|q|^2,
\end{aligned}
\end{equation}

we see that the structure condition (14.9) will imply (14.16) provided we choose

\[
\mubar=4\bigl(\mu(1+|\varphi|_1^2)+|\varphi|_2\bigr).
\]

Consequently, the estimate (14.15) will hold with $u$ replaced by $u-\varphi$ and $\mu$ replaced by $\mubar$. We can therefore assert the following boundary gradient estimate.

\begin{theorem}\label{gt14-boundary-gradient-14-1}\dcref{Theorem 14.1}\group{ch14-boundary-gradient}\source{gilbarg-trudinger:page-0349}
\textbf{Theorem 14.1.} \emph{Let $u\in C^2(\Omega)\cap C^1(\bar{\Omega})$ satisfy $\Q u=0$ in $\Omega$ and $u=\varphi$ on $\partial\Omega$. Suppose that $\Omega$ satisfies a uniform exterior sphere condition and $\varphi\in C^2(\bar{\Omega})$. Then if the structure condition (14.9) holds, we have}
\end{theorem}

\begin{equation}\tag{14.18}
|Du|\leq C \qquad \text{on }\partial\Omega
\end{equation}

\emph{where $C=C(n, M, \mu(M), \Phi, \delta)$, $M=\sup_{\Omega}|u|$, $\Phi=|\varphi|_{2;\Omega}$ and $\delta$ is the radius of the assumed exterior spheres.}

It is often convenient to write condition (14.9) in the form

\begin{equation}\tag{14.19}
p\Lambda,\, b = O(\mathcal{E}^0) \qquad \text{as } |p|\to\infty
\end{equation}

where the limit behavior with respect to $|p|$ is understood to be uniform in $\Omega\times(-N,N)$ for any $N>0$. In particular, if $\Q$ is uniformly elliptic in $\Omega\times(-N,N)\times\mathbb{R}^n$ for any $N>0$, that is $\Lambda=O(\lambda)$, and if also $b=O(\lambda|p|^2)$, then the structure condition (14.9) is fulfilled.

\section*{14.2 Convex Domains}

We consider in this section barrier constructions that are applicable to convex and uniformly convex domains. Suppose that $\Omega$ satisfies an exterior plane condition at a point $x_0\in\partial\Omega$, so that there exists a hyperplane $\Pcal$ with $x_0\in\Pcal\cap\bar{\Omega}=\Pcal\cap\partial\Omega$. Setting $d(x)=\operatorname{dist}(x,\Pcal)$ and $w=\psi(d)$ we then obtain, for any $u\in C^1(\bar{\Omega})\cap C^2(\Omega)$,

\begin{equation}\tag{14.20}
\bar{\Q}w=\psi' a^{ij}D_i d\, D_j d+b+\frac{\psi''}{(\psi')^2}\mathcal{E}^0
\end{equation}

\[
=b+\frac{\psi''}{(\psi')^2}\mathcal{E}^0
\]


% PDF page 350
\source{gilbarg-trudinger:page-0350}

Hence if $b = O(\mathscr{E})$, so that for some non-decreasing function $\mu$ we have

\begin{equation}
|b| \leq \mu(|z|)\,\mathscr{E} \qquad \text{for } |p| \geq \mu(|z|),
\tag{14.21}
\end{equation}

the barrier argument of the preceding section is applicable with $\nu = \mu(M)$, $M = \sup_{\Omega} |u|$.
Consequently we obtain an estimate for $Du(x_0)$, provided $Qu = 0$ in $\Omega$ and $u = 0$
on $\partial \Omega$. In order to extend this result to non-zero boundary values $\varphi$, we require
that $\Lambda D^2 \varphi$ and $b = O(\mathscr{F})$, that is,

\begin{equation}
\Lambda|D^2 \varphi| + |b| \leq \tilde{\mu}(|z|)\,\mathscr{F}
\qquad \text{for } |p - D\varphi| \geq \tilde{\mu}(|z|)
\tag{14.22}
\end{equation}

for some non-decreasing function $\tilde{\mu}$. We therefore have the following estimate.

\textbf{Theorem 14.2.} \emph{Let $u, \varphi \in C^2(\Omega)\cap C^1(\bar{\Omega})$ satisfy $Qu = 0$ in $\Omega$ and $u = \varphi$ on $\partial \Omega$
and suppose that $\Omega$ is convex. Then if the structure condition (14.22) holds, we have}

\begin{equation}
|Du| \leq C \qquad \text{on } \partial \Omega
\tag{14.23}
\end{equation}

\emph{where $C = C(n, M, \tilde{\mu}(M), |\varphi|_{1;\Omega})$, $M = \sup_{\Omega} |u|$.}

As in the preceding section, the structure condition (14.22) can be replaced in
the hypotheses of Theorem 14.2 by conditions that are independent of the boundary
values $\varphi$. In particular, either the condition

\begin{equation}
\Lambda = o(\mathscr{E}),\ b = O(\mathscr{E}) \qquad \text{as } |p| \to \infty,
\tag{14.24}
\end{equation}

or the condition

\begin{equation}
\Lambda,\, b = O(\lambda |p|^2) \qquad \text{as } |p| \to \infty
\tag{14.25}
\end{equation}

implies the validity of (14.22) for some function $\tilde{\mu}$ depending on $|\varphi|_{2;\Omega}$.
The first implication is a consequence of inequality (14.17); the second follows from the
inequality

\begin{equation}
\mathscr{F}(x, z, p, q) \geq \lambda(x, z, p)|p-q|^2, \qquad (x, z, p, q) \in \Omega \times \mathbb{R} \times \mathbb{R}^n \times \mathbb{R}^n.
\tag{14.26}
\end{equation}

We can therefore assert the following consequence of Theorem 14.2.

\begin{corollary}\label{gt14-boundary-gradient-14-3}\dcref{Corollary 14.3}\group{ch14-boundary-gradient}\source{gilbarg-trudinger:page-0350}
\textbf{Corollary 14.3.} \emph{Let $u \in C^2(\Omega)\cap C^1(\bar{\Omega})$ satisfy $Qu = 0$ in $\Omega$ and $u = \varphi$ on $\partial \Omega$.
Suppose that $\Omega$ is convex and $\varphi \in C^2(\bar{\Omega})$. Then, if either of the structure conditions
(14.24), (14.25) hold, we have}
\end{corollary}

\begin{equation}
|Du| \leq C \qquad \text{on } \partial \Omega,
\tag{14.27}
\end{equation}

\emph{where $C = C(n, M, \tilde{\mu}, |\varphi|_{2;\Omega})$.}

% PDF page 351
\source{gilbarg-trudinger:page-0351}

14.2. Convex Domains

Corollary 14.3 is in particular applicable to the minimal surface operator \(\M\)
given by

\[
(14.28)\qquad \M u=(1+|Du|^2)\,\Delta u-D_i u\,D_j u\,D_{ij}u.
\]

Here \(\lambda=1\), \(\Lambda=1+|p|^2\) (see Chapter 10, Examples (i) and (iii)), so that the
boundary gradient estimate is valid in convex domains for the equation \(\M u=0\).
This result which is already sharp in two dimensions, will be improved for higher
dimensions in the next section.

Let us next suppose that the domain \(\Omega\) satisfies an enclosing sphere condition at
a point \(x_0\in\partial\Omega\), that is there exists a ball \(B=B_R(y)\supset\Omega\) with \(x_0\in\partial B\). Setting
\(d(x)=\operatorname{dist}(x,\partial B)\) and \(w=\psi(d)\), we thus obtain, for any \(u\in C^1(\overline{\Omega})\cap C^2(\Omega)\),

\[
\overline{Q}w=\psi' a^{ij}D_{ij}d+b+\frac{\psi''}{(\psi')^2}\,\Ecal
\]

\[
(14.29)\qquad \le -\frac{\psi'}{R}\,\bigl(\Tcal-\Ecal^*\bigr)+b+\frac{\psi''}{(\psi')^2}\,\Ecal,
\]

where \(\Tcal(x,z,p)=\operatorname{trace}[a^{ij}(x,z,p)]=a^{ii}(x,z,p)\) and \(\Ecal^*=\Ecal/|p|^2\). The last relation
follows since

\[
D_{ij}d=-|x-y|^{-3}\bigl(|x-y|^2\delta_{ij}-(x_i-y_i)(x_j-y_j)\bigr).
\]

A boundary gradient estimate can now be derived from (14.29) if \(b\) is bounded in
terms of either \(\Tcal\) or \(\Ecal\). Indeed, let \(\varphi\in C^2(\overline{\Omega})\) and assume there exists a non-
decreasing function \(\bar\mu\) such that

\[
(14.30)\qquad |a^{ij}D_{ij}\varphi+b|\le \frac{1}{R}|p-D\varphi|\,\Tcal+\bar\mu\Ecal
\qquad\text{for } |p-D\varphi|\ge \bar\mu.
\]

The domain \(\Omega\) is said to be uniformly convex if it satisfies an enclosing sphere
condition at each boundary point with a ball of fixed radius \(R\). Our previous
barrier constructions then yield the following estimate.

\begin{theorem}\label{gt14-boundary-gradient-14-4}\dcref{Theorem 14.4}\group{ch14-boundary-gradient}\source{gilbarg-trudinger:page-0351}
Theorem 14.4.\quad Let \(u,\varphi\in C^2(\Omega)\cap C^1(\overline{\Omega})\) satisfy \(Qu=0\) in \(\Omega\) and \(u=\varphi\) on \(\partial\Omega\).
Suppose that \(\Omega\) is uniformly convex. Then if the structure condition (14.30) holds,
we have

\[
(14.31)\qquad |Du|\le C \quad \text{on }\partial\Omega
\]

where \(C=C(n,M,\bar\mu(M),R,|\varphi|_{1;\Omega})\).
\end{theorem}

It is clear that the structure condition (14.30) with \(\bar\mu\) depending on \(|\varphi|_{2;\Omega}\) is
implied by either the condition

\[
(14.32)\qquad b=o(\lambda|p|)+O(\lambda|p|^2) \quad \text{as } |p|\to\infty
\]


% PDF page 352
\source{gilbarg-trudinger:page-0352}




or the conditions

\begin{equation}
\Lambda = O(\lambda|p|^2),\ |b| \le \frac{|p|}{R}\,\Tcal + O(\lambda|p|^2) \qquad \text{as } |p| \to \infty.
\tag{14.33}
\end{equation}

Therefore we have the following consequence of Theorem 14.4.

\begin{corollary}\label{gt14-boundary-gradient-14-5}\dcref{Corollary 14.5}\group{ch14-boundary-gradient}\source{gilbarg-trudinger:page-0352}
\textbf{Corollary 14.5.} \emph{Let $u \in C^2(\Omega) \cap C^1(\bar{\Omega})$ satisfy $\Q u = 0$ in $\Omega$ and $u = \varphi$ on $\partial\Omega$.
Suppose that $\Omega$ is uniformly convex and $\varphi \in C^2(\bar{\Omega})$. Then, if either of the structure
conditions (14.32) or (14.33) hold, we have}

\begin{equation}
|Du| \le C \qquad \text{on } \partial\Omega
\tag{14.34}
\end{equation}

\emph{where $C = C(n, M, \mubar(M), |\varphi|_2, R)$.}
\end{corollary}

Corollary 14.5 is in particular applicable to the prescribed mean curvature
equation

\begin{equation}
\Mcal u = nH(x, u, Du)(1 + |Du|^2)^{3/2}.
\tag{14.35}
\end{equation}

Here $\Tcal = 1 + (n - 1)(1 + |p|^2)$ so that a boundary gradient estimate will hold for
solutions of (14.35) in uniformly convex domains provided the function $H$ satisfies

\begin{equation}
|H| \le \frac{n - 1}{nR} \qquad \text{for } |p| \ge \mu(|z|).
\tag{14.36}
\end{equation}

This result will be improved for dimensions higher than two in the next section.
The structure condition (14.32) is obviously satisfied when $b = 0$. In this case
Corollary 14.5 can be derived by means of linear barrier functions, since the
boundary manifold $(\partial\Omega,\varphi)$ will satisfy a bounded slope condition. Moreover when
$b = o(\lambda|p|)$ in (14.32) in Corollary 14.5, a barrier of the form $w = kd$, $d = \operatorname{dist}(x,\partial B)$
for sufficient large $k$ will suffice for the proof. It is also clear from the above proofs
that the results of this section will continue to be valid if the structure conditions
assumed in their hypotheses only hold for $x$ in some neighborhood of $\partial\Omega$.

\begin{remark}\label{gt14-boundary-gradient-14-r1}\dcref{Remark 14.1}\group{ch14-boundary-gradient}\source{gilbarg-trudinger:page-0352}
\textit{Remark.} \quad With slight modifications the considerations of this and the preceding
section can be combined, so that a boundary gradient estimate is determined only
by the \textit{local behavior} of the boundary $\partial\Omega$. It is convenient here to formulate such
a result for $C^2$ domains. Let $\kappa = \kappa(x_0)$ be the minimum of the principal curvatures
of $\partial\Omega$ at $x_0$ and let $\nu = \nu(x_0)$ denote the inner normal to $\partial\Omega$ at $x_0$. Suppose there
exists a non-decreasing function $\mubar$ such that for each $x_0 \in \partial\Omega$ there is an $\epsilon > 0$
for which

\begin{equation}
|a^{ij}D_{ij}\varphi + b| < \kappa\Tcal|p| + \mubar(|z|)\F
\qquad \forall(x, z, p) \in \Omega \times \mathbb{R} \times \mathbb{R}^n
\tag{14.37}
\end{equation}
\end{remark}


% PDF page 353
\source{gilbarg-trudinger:page-0353}

with

\[
\left|x-x_0\right|<\varepsilon,\qquad
\left|\frac{p-D\varphi}{\left|p-D\varphi\right|}\pm v\right|<\varepsilon
\quad\text{and}\quad
\left|p-D\varphi\right|\ge \mu.
\]

Then we have the estimate

\begin{equation}
\left|Du\right|\le C \qquad \text{on } \partial\Omega
\tag{14.38}
\end{equation}

where $C=C(n, M, \bar{\mu}(M), x)$. If $\kappa(x_0)\ge \kappa_0=\text{constant}$, for all $x_0\in\partial\Omega$ we can take the non-strict inequality in (14.37) provided $\kappa$ is replaced by $\kappa_0$. As an illustration of the possible type of behavior encompassed here, let us consider the equation

\begin{equation}
Qu=D_{11}u+\left(1+\left|D_2u\right|^N\right)D_{22}u=0,\qquad N\ge 0.
\tag{14.39}
\end{equation}

A boundary gradient estimate will then hold for a convex domain $\Omega$ and arbitrary $\varphi\in C^2(\bar{\Omega})$, provided $\kappa(x_0)>0$ whenever $v(x_0)\neq (\pm 1,0)$, that is provided the curvature of $\partial\Omega$ is positive except possibly when the tangent line is parallel to the $x_1$ axis.

\section*{14.3. Boundary Curvature Conditions}

So far in this chapter we have constructed barriers in terms of the distance function of an exterior surface of constant curvature (plane or sphere), the curvature of the latter being the significant factor in determining the structure condition imposed on the operator $Q$. By allowing more general exterior surfaces, our previous convexity conditions can be considerably relaxed in more than two dimensions. In the following we shall assume that $\partial\Omega\in C^2$ and use the boundary $\partial\Omega$ itself as an appropriate exterior surface. Setting $d(x)=\operatorname{dist}(x,\partial\Omega)$ we see from Lemma 14.16 that $d\in C^2(\Gamma)$ where $\Gamma=\{x\in\bar{\Omega}\mid d(x)<d_0\}$ for some $d_0>0$. Therefore if $w=\psi(d)$ where $\psi\in C^2[0,\infty)$ and $\psi'>0$ we have by formula (14.4), for any $u\in C^1(\bar{\Omega})\cap C^2(\Omega)$,

\begin{equation}
\bar{Q}w=a^{ij}(x,u(x),Dw)D_{ij}w+b(x,u(x),Dw)
\tag{14.40}
\end{equation}

\[
=\psi'a^{ij}D_{ij}d+b+\frac{\psi''}{(\psi')^2}\,\mathcal E.
\]

As a preliminary illustration of the general theory of this section, let us consider the special case of the minimal surface operator $\mathfrak{M}$; that is, we take

\begin{equation}
a^{ij}(x,z,p)=(1+|p|^2)\delta_{ij}-p_ip_j.
\tag{14.41}
\end{equation}

We then have

\begin{equation}
a^{ij}D_{ij}d=(1+|\psi'|^2)\Delta d-|\psi'|^2D_id\,D_jd\,D_{ij}d
\tag{14.42}
\end{equation}

\[
=(1+\psi'^2)\Delta d \qquad \text{since } |Dd|=1,\ D_i d D_{ij}d=0
\]

\[
\le -(n-1)(1+|\psi'|^2)H' \qquad \text{by Lemma 14.17,}
\]


% PDF page 354
\source{gilbarg-trudinger:page-0354}

where $H'$ is the mean curvature of $\partial\Omega$ at the point $y=y(x)$ on $\partial\Omega$ closest to $x$.
Hence, if $\partial\Omega$ has non-negative mean curvature everywhere, we obtain
\[
\Qbar w\le b+\frac{\psi''}{|\psi'|^2}\mathcal{E}
\]
as in the convex case of the previous section. A boundary gradient estimate
then follows as in the preceding section for arbitrary $C^2(\bar\Omega)$ boundary values if
$b=O(|p|^2)$. We shall show in the next section that this result is sharp for the
minimal surface equation $\M u=0$. Using relations (14.40) and (14.42) we can also
conclude a corresponding sharp result for the equation of prescribed mean
curvature (14.35). However, let us first return to the general situation. We assume
that the coefficients of $Q$ are decomposed in such a way that for $p\neq 0$ we have
\begin{equation}
\tag{14.43}
\begin{aligned}
 a^{ij} &= \lambda a^{ij}_{\infty}+a^{ij}_0, \qquad i,j=1,\ldots,n, \\
 b &= |p|\lambda b_{\infty}+b_0
\end{aligned}
\end{equation}
where
\[
 a^{ij}_{\infty}(x,z,p)=a^{ij}_{\infty}(x,z,p/|p|), \qquad b_{\infty}(x,z,p)=b_{\infty}(x,z,p/|p|),
\]
\[
 a^{ij}_{\infty}(x,z,\xi)\,\xi_i\xi_j\ge 0 \qquad \text{for all } x\in\Omega,\ |\xi|=1,\ \xi\in\R^n
\]
and $b_{\infty}$ is non-increasing in $z$. For example, in the case of the minimal surface
operator $\M$ we can take
\[
 a^{ij}_{\infty}=\delta_{ij}-p_i p_j/|p|^2, \qquad a^{ij}_0=p_i p_j/|p|^2.
\]
Using the matrix $[a^{ij}_{\infty}]$, we introduce a generalized notion of mean curvature as
follows. Namely, let $y$ be a point of $\partial\Omega$ and let $\nu$ denote the unit inner normal to
$\partial\Omega$ at $y$, $\kappa_1,\ldots,\kappa_{n-1}$ the principal curvatures of $\partial\Omega$ at $y$ and
$a_1,\ldots,a_n$ the diagonal elements of the matrix $[a^{ij}_{\infty}]$ with respect to a corresponding principal coordinate
system at $y$. We then define
\begin{equation}
\tag{14.44}
\xcal^{\pm}(y)=\sum_{i=1}^{n-1} a_i(y,\pm\nu)\kappa_i.
\end{equation}
Since $a_i\ge 0$, $i=1,\ldots,n$, the quantities $\xcal^{\pm}$ are a weighted average of the curvatures
of $\partial\Omega$ at $y$. Furthermore, in the special case of the minimal surface operator $\M$,
we have $a_i=1$, $i=1,\ldots,n-1$, $a_n=0$ and hence
\[
\xcal^+(y)=\xcal^-(y)=\sum_{i=1}^{n-1}\kappa_i=(n-1)H'(y)
\]


% PDF page 355
\source{gilbarg-trudinger:page-0355}


where $H'(y)$ denotes the mean curvature of $\partial\Omega$ at $y$. Through Lemma 14.17 we see
that the curvatures $\mathcal{K}^{\pm}$ are connected with the distance function $d$ by the formula

\begin{equation}
\tag{14.45}
\mathcal{K}^{\pm}(y)=-a^{ij}_{\infty}(y,\pm Dd(y))D_{i}d(y).
\end{equation}

In order to generalize our earlier result for the minimal surface equation let us
suppose that the inequality

\begin{equation}
\tag{14.46}
\mathcal{K}^{+}\ge b_{\infty}(y,u,v)
\end{equation}

holds at each point $y\in\partial\Omega$. In addition we assume that the functions $a^{ij}_{\infty}$, $b_{\infty}\in$
$C^{1}(\Gamma\times\R\times\R^{n})$ and that $\cQ$ satisfies the structure condition

\begin{equation}
\tag{14.47}
\Lambda,\, |p|\,|a^{ij}_{0}|,\, b_{0}=O(\cE) \qquad \text{as } |p|\to\infty,\ i,j=1,\ldots,n,
\end{equation}

so that for some non-decreasing function $\mu$ we have

\begin{equation}
\tag{14.48}
\Lambda+|p|\sum |a^{ij}_{0}|+|b_{0}|\le \mu(|z|)\cE \qquad \text{for } |p|\ge \mu(|z|).
\end{equation}

As previously we shall assume initially that the function $u$ vanishes on $\partial\Omega$. Then
taking $w$ as before we have

\begin{align*}
\bar{Q}^{\,w}
&=a^{ij}(x,u(x),Dw)D_{ij}w+|Dw|\Lambda(x,u(x),Dw)b_{\infty}(x,w,Dw) \\
&\qquad +b_{0}(x,u(x),Dw) \\
&=\psi'\Lambda(a^{ij}_{\infty}D_{ij}d+b_{\infty})+\psi'a^{ij}_{0}D_{ij}d+b_{0}+\frac{\psi''}{|\psi'|^{2}}\cE
\end{align*}

by (14.4) and (14.43), where

\begin{align*}
a^{ij}_{\infty}D_{ij}d+b_{\infty}
&=a^{ij}_{\infty}(x,v)D_{ij}d(x)+b_{\infty}(x,w,v) \\
&\le a^{ij}_{\infty}(x,v)D_{ij}d(y)+b_{\infty}(x,0,v) \qquad \text{by Lemma 14.17} \\
&\le \bigl(a^{ij}_{\infty}(x,v)-a^{ij}_{\infty}(y,v)\bigr)D_{ij}d(y)+b_{\infty}(x,0,v)-b_{\infty}(y,0,v)
\end{align*}

by (14.45) and (14.46)

\[
\le Kd.
\]

Here $y=y(x)$ is the point on $\partial\Omega$ closest to $x$, $v=Dd(y)=Dd(x)$ is the inner unit
normal to $\partial\Omega$ at $y$ and the constant $K$ is given by

\begin{equation}
\tag{14.49}
K=\sup_{x\in\Gamma}\frac{\bigl|(a^{ij}_{\infty}(x,v)-a^{ij}_{\infty}(y,v))D_{ij}d(y)+b_{\infty}(x,u(y),v)-b_{\infty}(y,u(y),v)\bigr|}{|x-y|}.
\end{equation}


% PDF page 356
\source{gilbarg-trudinger:page-0356}



Hence by (14.48) we obtain

\[
\bar{\mathcal Q}w \le K\psi' d\mu\delta + \bigl(1+\sup_{\Gamma}|D^2 d|\bigr)\mu\delta
  + \frac{\psi''}{|\psi'|^2}\,\delta
\]
\[
\le \left(\nu+\frac{\psi''}{|\psi'|^2}\right)\delta
\]

provided $\psi'd\le 1$ and $\psi'\ge \mu$ where $\nu=(K+1+\sup_{\Gamma}|D^2 d|)\mu$, $\mu=\mu(M)$ and
$M=\sup_{\Gamma}|u|$. Consequently by choosing $\psi$ by formula (14.11) and taking $k$ large enough to guarantee
$a\le d_0$, the function $w^+=w$ will be an upper barrier at every point of $\partial\Omega$ for the operator
$\bar{\mathcal Q}$ and the function $u$. A lower barrier is similarly constructed provided instead of (14.46) the
inequality

\[
(14.46)'\qquad \mathfrak X^- \ge -\,b_\infty(y,u(y),-v)
\]

holds at each point $y\in\partial\Omega$. Hence, if both (14.46) and (14.46)' hold and $\mathcal Q u=0$ in $\Omega$,
then $u$ will satisfy the estimate (14.3) at each $x_0\in\partial\Omega$. In order to extend these results to
non-zero boundary values $\varphi$, we require that $\Lambda$, $|p|a_0^{ij}$ and $b_0=O(\mathfrak F)$, that is,

\[
(14.50)\qquad \Lambda + |p| \sum |a_0^{ij}| + |b_0|\le \bar\mu(|z|)\mathfrak F
 \qquad \text{for } |p-D\varphi|\le \bar\mu(|z|)
\]

for some non-decreasing function $\bar\mu$. For the transformed operator $\tilde{\mathcal Q}$ given by (14.5) we
can take $\tilde a_\infty^{ij}=a_\infty^{ij}$, $\tilde b_\infty(x,z+\varphi,p)=b_\infty(x,z,p)$ so that conditions
(14.46) and (14.46)' are unchanged. We therefore have the following estimate.

\medskip

\begin{theorem}\label{gt14-boundary-gradient-14-6}\dcref{Theorem 14.6}\group{ch14-boundary-gradient}\source{gilbarg-trudinger:page-0355}
\noindent\textbf{Theorem 14.6.} \textit{Let $u\in C^2(\Omega)\cap C^1(\bar\Omega)$, $\varphi\in C^2(\bar\Omega)$
satisfy $\mathcal Q u=0$ in $\Omega$ and $u=\varphi$ on $\partial\Omega$. Suppose that $\mathcal Q$ satisfies the
structure conditions (14.43), (14.50) and that at each point $y\in\partial\Omega$ the inequalities}

\[
(14.51)\qquad \mathfrak X^+ \ge b_\infty(y,\varphi(y),v),\qquad
\mathfrak X^- \ge -b_\infty(y,\varphi(y),-v)
\]

\textit{hold. Then we have the estimate}

\[
(14.52)\qquad |Du|\le C \quad \text{on }\partial\Omega
\]

\textit{where $C=C(n,M,\bar\mu(M),\Omega,K,|\varphi|_{2;\Omega})$, $M=\sup_{\Omega}|u|$ and $K$ is given by
(14.49) with $u$ replaced by $\varphi$.}
\end{theorem}

\medskip

As in the case of convex $\Omega$ treated in Section 14.2, the structure condition (14.50) can be replaced in the
hypotheses of the above theorem by conditions that are independent of the boundary values $\varphi$. In particular
either the condition

\[
(14.53)\qquad \Lambda=o(\mathfrak F),\qquad |p|a_0^{ij},\, b_0=O(\mathfrak F)\quad \text{as }|p|\to\infty,
\]


% PDF page 357
\source{gilbarg-trudinger:page-0357}


or the condition

\[
(14.54)\qquad \Lambda, |p|a^{ij}_0, b_0 = O(\lambda |p|^2)\quad \text{as } |p| \to \infty,
\]

implies the validity of (14.50) for some function $\bar{\mu}$ depending on
$|\varphi|_{2;\Omega}$. We can therefore assert the following consequence of
Theorem 14.6.

\medskip

\begin{corollary}\label{gt14-boundary-gradient-14-7}\dcref{Corollary 14.7}\group{ch14-boundary-gradient}\source{gilbarg-trudinger:page-0357}
\noindent\textbf{Corollary 14.7.} \textit{Let $u \in C^2(\Omega) \cap C^1(\bar{\Omega})$,
$\varphi \in C^2(\bar{\Omega})$ satisfy $Qu = 0$ in $\Omega$ and $u = \varphi$
on $\partial\Omega$. Suppose that, in addition to (14.43), either of the
structure conditions (14.53), (14.54) hold and that the inequalities (14.51)
are satisfied on $\partial\Omega$. Then we have the estimate}

\[
(14.55)\qquad |Du| \le C \quad \text{on } \partial\Omega,
\]

\noindent\textit{where $C = C(n, M, \bar{\mu}, \Omega, K, |\varphi|_{2;\Omega})$.}
\end{corollary}

\medskip

\noindent The application of Corollary 14.7 to the equation of prescribed mean
curvature,

\[
(14.56)\qquad \mathfrak M u = nH(x,u)(1+|Du|^2)^{3/2},
\]

\noindent where $H \in C^1(\bar{\Omega} \times \mathbb R)$ and $D_z H \ge 0$, now
yields the following result.

\medskip

\begin{corollary}\label{gt14-boundary-gradient-14-8}\dcref{Corollary 14.8}\group{ch14-boundary-gradient}\source{gilbarg-trudinger:page-0357}
\noindent\textbf{Corollary 14.8.} \textit{Let $u$ be a $C^2(\Omega) \cap C^1(\bar{\Omega})$
solution of equation (14.35) in $\Omega$ with $u = \varphi$ on $\partial\Omega$
where $\varphi \in C^2(\bar{\Omega})$. Suppose that the mean curvature $H'$ of
$\partial\Omega$ is such that}

\[
(14.57)\qquad H'(y) \ge \frac{n}{n-1}|H(y,\varphi(y))| \qquad \forall\, y \in \partial\Omega.
\]

\noindent\textit{Then we have the estimate}

\[
(14.58)\qquad |Du| \le C \quad \text{on } \partial\Omega,
\]

\noindent\textit{where $C = C(n, M, \Omega, H_1, |\varphi|_{2;\Omega})$, $M = \sup_{\Omega}|u|$,
$H_1 = \sup_{\Omega \times (-M,M)} (|H| + |DH|)$.}
\end{corollary}

\medskip

The sharpness of Corollary 14.8 will be demonstrated in the next section. Note
that the result could have been derived more directly by proceeding along the
lines of our earlier treatment of the minimal surface equation, $\mathfrak M u = 0$.
So far the results in this section have necessitated some control over the
behavior of the maximum eigenvalue $\lambda$ with respect to $\FF$, $\EE$ or
$\lambda$. We move on now to consider a situation where such controls are not
imposed but where, to compensate, the inequalities (14.51) must hold in the
strict sense everywhere on the boundary $\partial\Omega$. The forerunner of this
situation is the case in Corollary 14.5 where the structure condition (14.32)
holds. To proceed further we assume that in the decomposition (14.43) the
coefficients $a^{ij}_\infty, b_\infty$ are continuous on $\partial\Omega \times
\mathbb R \times \mathbb R^n$ and that the coefficients

\[
(14.59)\qquad a^{ij}_0 = o(\Lambda),\ b_0 = o(|p|\Lambda)\quad \text{as } |p| \to \infty.
\]

% PDF page 358
\source{gilbarg-trudinger:page-0358}

Let $u \in C^0(\overline{\Omega}) \cap C^2(\Omega)$ satisfy $Qu = 0$ in $\Omega$ and
$u = \varphi$ on $\partial\Omega$ where $\varphi \in C^2(\overline{\Omega})$, and
suppose that the strict inequalities

\begin{equation}\tag{14.60}
\mathcal{K}^+ > b_\infty(y,\varphi(y),v), \qquad
\mathcal{K}^- > -b_\infty(y,\varphi(y),-v)
\end{equation}

hold everywhere on $\partial\Omega$. Assume initially $\varphi \equiv 0$ so that setting
$w = kd$ for some positive constant $k$, we have

\[
\overline{Q}w = a^{ij}(x,u(x),Dw)D_{ij}w + |Dw|\Lambda(x,u(x),Dw)b_\infty(x,w,Dw)
+ b_0(x,u(x),Dw)
\]
\[
= k\Lambda(a^{ij}_\infty(x,Dd)D_{ij}d + b_\infty(x,w,Dd)) + ka^{ij}_0D_{ij}d + b_0
\]
\[
= \Lambda\{k(a^{ij}_\infty(x,Dd)D_{ij}d + b_\infty(x,w,Dd)) + o(k)\}
\]

by (14.59). Now by Lemma 14.16, the relation (14.45) and the fact that $b_\infty$ is non-increasing in $z$, there exist positive constants $\chi$ and $a \le d_0$ such that

\[
a^{ij}_\infty(x,Dd)D_{ij}d + b_\infty(x,w,Dd) \le -\chi
\]

in the neighborhood $\mathcal{N} = \{x \in \overline{\Omega}\mid d(x) < a\}$. Hence we obtain

\[
\overline{Q}w \le \Lambda(-k\chi + o(k)) < 0
\]

for sufficiently large $k$.

Consequently, by choosing $k$ large enough so that also $ka \ge \sup_{\Omega}|u|$, the function
$w^+ = w$ will be an upper barrier at every point of $\partial\Omega$ for the operator $\overline{Q}$ and the function $u$. Similarly the function $w^- = -w$ will be a corresponding lower barrier and we obtain $|Du| \le k$ on $\partial\Omega$ for $u \in C^1(\overline{\Omega})$. This result extends automatically to non-zero boundary values $\varphi$ by replacement of $u$ by $u - \varphi$ and use of (14.5). Thus we have proved the following estimate.

\begin{theorem}\label{gt14-boundary-gradient-14-9}\dcref{Theorem 14.9}\group{ch14-boundary-gradient}\source{gilbarg-trudinger:page-0358}
\noindent\textbf{Theorem 14.9.} \textit{Let $u \in C^2(\Omega) \cap C^1(\overline{\Omega})$, $\varphi \in C^2(\overline{\Omega})$ satisfy $Qu = 0$ in $\Omega$ and $u = \varphi$ on $\partial\Omega$. Suppose that $Q$ satisfies the structure conditions (14.43), (14.59) and that the inequalities (14.60) hold at each point $y \in \partial\Omega$. Then we have the estimate}

\begin{equation}\tag{14.61}
|Du| \le C \qquad \text{on } \partial\Omega
\end{equation}

\[
\text{where } C = C(n,M,a^{ij}_\infty,a^{ij}_0,b_\infty,b_0,\Omega,|\varphi|_{2;\Omega})
\text{ and } M = \sup_{\Omega}|u|.
\]
\end{theorem}

The dependence of the constant $C$ in Theorem 14.9 on the coefficients $a^{ij}_0$, $b_0$
arises through the structure conditions (14.59), and the dependence on the coefficients
$a^{ij}_\infty$, $b_\infty$ through their moduli of continuity on $\partial\Omega \times \mathbb{R} \times \mathbb{R}^n$.

To conclude this section, let us note the relationship of the results here to those
of the preceding sections. If we choose $a^{ij}_\infty = b_\infty = 0$ in the decomposition (14.43),
then Theorem 14.6 reduces to Theorem 14.1 for $\partial\Omega \in C^2$. Next although the results
in Section 14.2 on convex domains are not special cases of Theorems 14.6 and 14.9
they are encompassed by slight variants of these theorems which are treated in
Problems 14.2 and 14.3.


% PDF page 359
\source{gilbarg-trudinger:page-0359}



\section*{14.4. Non-Existence Results}

We present here some non-existence results which show that many of the conditions in the theorems of the preceding sections cannot be significantly relaxed. Since the class of equations considered in this section include equations for which Steps I and III of the existence procedure, described in Chapter 11, are readily established, it follows that the non-solvability of the Dirichlet problem for these cases is due to the lack of a boundary gradient estimate. Indeed the non-existence of such an estimate for these equations can be demonstrated directly by techniques analogous to those used below.

The main tool for our treatment of non-existence results is the following variant of the comparison principle, Theorem 10.1.

\begin{theorem}[14.10]\label{gt14-boundary-gradient-14-10}\dcref{Theorem 14.10}\group{ch14-boundary-gradient}\source{gilbarg-trudinger:page-0359}
Let $\Omega$ be a bounded domain in $\mathbb{R}^n$ and $\Gamma$ a relatively open $C^1$ portion of $\partial\Omega$. Then if $Q$ is an elliptic operator of the form (14.2) and $u \in C^0(\overline{\Omega}) \cap C^2(\Omega \cup \Gamma)$, $v \in C^0(\overline{\Omega}) \cap C^2(\Omega)$ satisfy $Qu \ge Qv$ in $\Omega$, $u \le v$ on $\partial\Omega - \Gamma$, $\partial v/\partial \nu = -\infty$ on $\Gamma$, it follows that $u \le v$ in $\Omega$.
\end{theorem}

\begin{proof}
By Theorem 10.1, we have
\[
\sup_{\Omega}(u-v)\le \sup_{\Gamma}(u-v)^+.
\]
Since
\[
\frac{\partial}{\partial \nu}(u-v)=\frac{\partial u}{\partial \nu}-\frac{\partial v}{\partial \nu}=\infty
\]
on $\Gamma$, the function $u-v$ cannot achieve a maximum value on $\Gamma$. Hence $u\le v$ in $\Omega$. $\square$
\end{proof}

In order to apply Theorem 14.10, we let $y\in \partial\Omega$ be fixed, $\delta = \operatorname{diam}\Omega, 0<a<\delta/2$, and consider the function $w$ defined by

\[
w(x)=m+\psi(r), \qquad r=|x-y|,
\]

where $m \in \mathbb{R}$ and $\psi \in C^2(a,\delta)$ is such that $\psi(\delta)=0$, $\psi' \le 0$, $\psi'(a)=-\infty$. Using (14.8), we obtain for $u \in C^2(\Omega)$, $r>a$,

\begin{equation}
\bar{Q}w=a^{ij}(x,u(x),Dw)D_{ij}w+b(x,u(x),Dw)
=\frac{\psi'}{r}\,(\mathcal{T}-\mathcal{E}^*)+\frac{\psi''}{(\psi')^2}\,\mathcal{E}+b,
\tag{14.62}
\end{equation}

the arguments of $\mathcal{T}=\operatorname{trace}[a^{ij}]$, $\mathcal{E}^*=\mathcal{E}/|p|^2$, $\mathcal{E}$ and $b being x, u(x) and Dw$. We now wish to choose the function $\psi$ in such a way that $\bar{Q}w<0$ in the domain $\bar{\Omega}=\{x\in \Omega \mid r>a,\ |u(x)|>M\}$ for some constant $M$. If this is done and $Qu=0$ in $\Omega$ we then have by Theorem 14.10

\begin{equation}
\sup_{\Omega-B_a(y)} u \le M+m+\psi(a)
\tag{14.63}
\end{equation}


% PDF page 360
\source{gilbarg-trudinger:page-0360}




348\hfill 14.\ Boundary Gradient Estimates

where

\[
m=\sup_{\partial\Omega-B_a(y)}u^+.
\]

The estimate (14.63) can be regarded as a preliminary stage in the establishment of non-existence results.

We shall consider two different cases.

(i) First, let us suppose that

\begin{equation}
b\leq -|p|^\theta \Ecal
\tag{14.64}
\end{equation}

for $x\in\Omega$, $|z|\geq M$, $|p|\geq L$ where $M$, $L$ and $\theta$ are positive constants. Then setting

\begin{equation}
\psi(r)=K[(\delta-a)^\beta-(r-a)^\beta]
\tag{14.65}
\end{equation}

where $\beta=\theta/(1+\theta)$ and $K\in\mathbb{R}$ we obtain, for sufficiently large $K$, $\Qbar w<0$ in $\Omegabar$. Hence the estimate (14.63) holds in this case.

(ii) Next, let us suppose that

\begin{equation}
b\leq 0,
\qquad
\Ecal\leq \mu(\mathcal{T}-\Ecal^\ast)|p|^{1-\theta}
\tag{14.66}
\end{equation}

for $x\in\Omega$, $|z|\geq M$, $|p|\geq 0$, where $\mu$, $\theta$ and $M$ are positive constants. Then we have, in $\Omegabar$,

\[
\Qbar w\leq \left(-\frac{|\psi'|^\theta}{\mu r}+\frac{\psi''}{(\psi')^2}\right)\Ecal<0
\]

for the specific choice

\begin{equation}
\psi(r)=\left(\frac{\mu}{\beta}\right)^\beta\int_r^\delta \left(\log\frac{t}{a}\right)^{-\beta}\,dt
\tag{14.67}
\end{equation}

where $\beta=1/(1+\theta)$. Hence again the estimate (14.63) follows. Note that the function $\psi$ given by (14.67) satisfies

\[
\lim_{a\to 0}\psi(a)=0.
\]

Now let us assume that the domain $\Omega$ satisfies an interior sphere condition at the point $y$ so that there exists a ball $B=B_R(x_0)\subset\Omega$ with $y\in\bar{B}\cap\partial\Omega$. We then consider


% PDF page 361
\source{gilbarg-trudinger:page-0361}


the function $w^{*}$ defined by

\begin{equation}
w^{*}(x)=m^{*}+\chi(r), \qquad r=|x-x_{0}|
\tag{14.68}
\end{equation}

where $m^{*}\in\R$ and $\chi\in C^{2}(0, R-\varepsilon)$, $0<\varepsilon<R$, is such that $\chi(0)=0$, $\chi'\geq 0$ and $\chi'(R-\varepsilon)=\infty$. By (14.62) we have, for $r>0$,

\begin{equation}
\bar{Q}w^{*}=\frac{\chi'}{r}\left(\mathcal{T}-\mathcal{E}^{*}\right)+\frac{\chi''}{(\chi')^{2}}\mathcal{E}^{*}+b.
\tag{14.69}
\end{equation}

Let us impose now a stronger condition than (14.64), namely that

\begin{equation}
b+\frac{|p|}{R'}\,\mathcal{T}\leq -|p|^{\theta}\mathcal{E}^{*}, \qquad 0<R'<R,
\tag{14.70}
\end{equation}

for $x\in\Omega$, $|z|\geq M$, $|p|\geq L$. Then setting

\begin{equation}
\chi(r)=K\{(R-\varepsilon)^{\beta}-(R-\varepsilon-r)^{\beta}\}, \qquad \beta=\frac{\theta}{1+\theta}.
\tag{14.71}
\end{equation}

we obtain, for sufficiently large $K$, $\bar{Q}w^{*}<0$ in the domain $\tilde{\Omega}=\{x\in\Omega \mid |x-y|<R,\ R'<|x-x_{0}|<R-\varepsilon,\ |u(x)|>M\}$. Hence by Theorem 14.10 we have

\begin{equation}
\sup_{\tilde{\Omega}} u\leq M+m^{*}+\chi(R-\varepsilon)
\tag{14.72}
\end{equation}

where

\[
m^{*}=\sup_{|x-y|=R} u.
\]

Combining the estimates (14.72) and (14.63) with $a=R-R'$, and letting $\varepsilon$ tend to zero, we therefore obtain the estimate

\begin{equation}
u(y)\leq 2M+m+K[(\delta-R)^{\beta}+R^{\beta}]
\tag{14.73}
\end{equation}

where $K$ depends on $\theta$ and $L$ and

\[
m=\sup_{\partial\Omega - B_{R}(y)} u.
\]

The estimate (14.73) shows that the boundary values of the function $u$ cannot be specified arbitrarily on $\partial\Omega$. Since the above argument can be repeated with $u$ replaced by $-u$, the structure condition (14.70) can be relaxed to

\begin{equation}
|b|\geq \frac{|p|}{R'}\,\mathcal{T}+|p|^{\theta}\mathcal{E}^{*}
\tag{14.74}
\end{equation}

% PDF page 362
\source{gilbarg-trudinger:page-0362}

for $x\in \Omega$, $|z|\ge M$, $|p|\ge L$. We have thus proved the following non-existence
result.

\begin{theorem}\label{gt14-boundary-gradient-14-11}\dcref{Theorem 14.11}\group{ch14-boundary-gradient}\source{gilbarg-trudinger:page-0362}
\textbf{Theorem 14.11.} \textit{Let $\Omega$ be a bounded domain in $\mathbb{R}^n$ and suppose that the operator $\Qop$
satisfies the structure condition (14.74) where $R$ is the radius of the largest ball
contained in $\Omega$. Then there exists a function $\varphi\in C^{\infty}(\overline{\Omega})$ such that the Dirichlet
problem $\Qop u=0$ in $\Omega$, $u=\varphi$ on $\partial\Omega$ is not solvable.}
\end{theorem}

Theorem 14.11 implies that both Theorems 14.1 and 14.4 are sharp in the sense
that the quantities $\E$, $\F$ in the structure conditions (14.9), (14.30) cannot be
replaced by $|p|^{\theta}\E$, $|p|^{\theta}\F$ for some $\theta>0$, (even if the operator $\Qop$ is the Laplacian!).
Furthermore, in Corollary 14.5, we can neither replace (14.32) by the condition
$b=O(\Lambda/|p|)$ as $|p|\to\infty$, nor in the second inequality in (14.33), replace $\lambda|p|^2$ by
$\lambda|p|^{2+\theta}$ for some $\theta>0$.

We shall use case (ii) above to demonstrate the need for the geometric restric-
tions in Section 14.3. Let us assume that the decomposition (14.43) is valid with
$b_x$ independent of $z$ and $a_{\alpha\beta}^{ij},\, b_x$ continuous on $\partial\Omega\times\mathbb{R}\times\mathbb{R}^n$ and satisfying the
structure conditions (14.59). In addition we impose on the operator $\Qop$ the structure
conditions

\[
\begin{array}{rcl}
(14.75) & b\le 0 \\
        & \E\le \mu \Lambda |p|^{1-\theta}
\end{array}
\]

for $x\in \Omega$, $|z|\ge M$, $|p|\ge 0$, where $\mu$, $\theta$ and $M$ are positive constants. It is easy to
show that the conditions (14.43), (14.59), (14.75) imply that condition (14.66)
holds for $x\in \Omega$, $M\le |z|\le \overline{M}$, $p\in \mathbb{R}^n$ for some constant $\overline{M}$ and a possibly different
constant $\mu$ than that in (14.75). Moreover, since (14.75) also implies, by the classical
maximum principle, Theorem 3.1, that

\[
(14.76)\qquad \sup_{\Omega} u \le M + \sup_{\partial\Omega} u,
\]

we can assume below that (14.66) is applicable in $\Omega^{+}=\{x\in \Omega\,|\,u(x)>0\}$ and that
also in $\Omega^{+}$ the quantities in (14.59) are bounded in terms of $\sup_{\partial\Omega}u$. Let us now
suppose that $\partial\Omega\in C^2$ and that

\[
(14.77)\qquad \Xminus(y)<-b_{\alpha\beta}(y,-\nu(y))
\]

where $\Xminus$ is given by (14.44) and $\nu$ denotes the unit inner normal to $\partial\Omega$ at $y$. The
argument which follows is analogous to the proof of Theorem 14.11, the interior
sphere at $y$ being replaced by another quadric surface. In fact, condition (14.77)
implies that for sufficiently small $a$, there exists a quadric surface $\Sscr$ tangent to $\partial\Omega$
at $y$ such that (i) $\Sscr$ has a unique parallel projection onto the tangent plane at $y$,
(ii) $\Sscr\cap B_a(y)\subset \overline{\Omega}$ and (iii) the curvature $\Xscr^{-}$ corresponding to $\Sscr$, $\Xscr_{\varphi}$, satisfies

\[
(14.78)\qquad \Xscr_{\varphi}(y)\le \Xminus(y)+\eta
\]


% PDF page 363
\source{gilbarg-trudinger:page-0363}

for some $\eta > 0$. We now consider the function $w^*$ defined by

\begin{equation}
w^*(x)=m^*+\chi(d), \qquad d=\operatorname{dist}(x,\mathcal{S}),
\tag{14.79}
\end{equation}

where $m^* \in \mathbb{R}$ and $\chi \in C^2(\epsilon,a),\, 0 < \epsilon < a$, is such that $\chi(2a)=0$, $\chi' \leq 0$, $\chi'(\epsilon)=-\infty$. By (14.40) we have, in the domain $\tilde{\Omega}=\{x \in \Omega \mid |x-y|<a,\, \epsilon<d<a,\, u(x)>M\}$,

\[
\overline{Q}w^*=\chi' \Lambda\!\bigl(a_{\infty}^{ij} D_{ij} d + b_{\infty}\bigr)+\chi' a_0^{ij} D_{ij} d + b_0 + \frac{\chi''}{(\chi')^2}\mathcal{E}
\]
\[
\leq \chi'\Lambda(\eta+o(1))+\frac{\chi''}{(\chi')^2}\mathcal{E}\qquad \text{by (14.59) and (14.78),}
\]
\[
\leq \chi'\Lambda(\eta+o(1))+\mu\chi'|x|^{-2-\theta}\qquad \text{by (14.75).}
\]

Setting

\begin{equation}
\chi(d)=K\bigl[(2a-\epsilon)^\beta-(d-\epsilon)^\beta\bigr],
\tag{14.80}
\end{equation}

where $\beta=\theta/(1+\theta)$, we then obtain, for sufficiently large $K$,

\[
\overline{Q}w^*<0
\]

in $\tilde{\Omega}$. Hence by Theorem 14.10, we have

\begin{equation}
\sup_{\tilde{\Omega}} u \leq M+m^*+\chi(\epsilon)
\leq M+m^*+K(2a)^\beta
\tag{14.81}
\end{equation}

where

\[
m^*=\sup_{|x-y|=a} u.
\]

Combining the estimates (14.81), (14.63) and letting $\epsilon$ tend to zero, we therefore obtain the estimate

\begin{equation}
u(y)\leq 2M+m+\psi(a)+K(2a)^\beta
\tag{14.82}
\end{equation}

where $\psi$ is given by (14.67) and

\[
m=\sup_{\partial\Omega-B_a(y)} u.
\]

Again, the estimate (14.82) shows that $u$ cannot be prescribed arbitrarily on $\partial\Omega$.
The sharpness of Theorem 14.6 is thus evidenced by the following non-existence result.

% PDF page 364
\source{gilbarg-trudinger:page-0364}


provided $b$ is replaced by $-b$ in the structure conditions (14.75). Also if $M=0$ in
(14.76) we can choose $\sup_{\partial\Omega}|\varphi|$ to be arbitrarily small. Specializing to the equation
of prescribed mean curvature (14.56) with $H(x, z)\equiv H(x)$, we therefore have

\begin{corollary}\label{gt14-boundary-gradient-14-13}\dcref{Corollary 14.13}\group{ch14-boundary-gradient}\source{gilbarg-trudinger:page-0364}
\textbf{Corollary 14.13.} \ Let $\Omega$ be a bounded $C^2$ domain in $\R^n$ and suppose that at some
point $y\in\partial\Omega$ the mean curvature $H'$ of $\partial\Omega$ is such that
\begin{equation}
H'(y)<\frac{n}{n-1}|H(y)|
\tag{14.85}
\end{equation}
where $H\in C^0(\overline{\Omega})$ is either non-positive or non-negative in $\Omega$. Then, for any $\varepsilon>0$,
there exists a function $\varphi\in C^\infty(\overline{\Omega})$ with $\sup|\varphi|\leq\varepsilon$, such that the Dirichlet problem,
\[
\Qop u=0 \text{ in } \Omega,\qquad u=\varphi \text{ on } \partial\Omega,
\]
is not solvable.
\end{corollary}

Combining Corollaries 14.8 and 14.13 with our treatment of the existence
procedure for the special case (11.7) in Section 11.3, we obtain the following sharp
criterion of Jenkins and Serrin for solvability of the minimal surface equation
[JS 2].

\begin{theorem}\label{gt14-boundary-gradient-14-14}\dcref{Theorem 14.14}\group{ch14-boundary-gradient}\source{gilbarg-trudinger:page-0364}
\textbf{Theorem 14.14.} \ Let $\Omega$ be a bounded $C^{2,\gamma}$ domain in $\R^n$, $0<\gamma<1$. Then the Dirichlet
problem $\Qop u=0$ in $\Omega$, $u=\varphi$ on $\partial\Omega$, is solvable for arbitrary $\varphi\in C^{2,\gamma}(\overline{\Omega})$ if and only
if the mean curvature $H'$ of $\partial\Omega$ is non-negative at every point of $\partial\Omega$.
\end{theorem}

The prescribed mean curvature equation will be studied further in Chapter 16.
We note here that the restriction that $b_\infty$ is independent of $z$ in Theorem 14.12 can
be removed and consequently that condition (14.85) in Corollary 14.13 can be
replaced by

\begin{equation}
H'(y)<\sup_{z\in\R}\frac{n}{n-1}|H(y,z)|;
\tag{14.86}
\end{equation}

(see Problem 14.5). It is also worth comparing the results of this section with the
existence theorems of Section 15.5 in the next chapter.

% PDF page 365
\source{gilbarg-trudinger:page-0365}


14.5. Continuity Estimates

The barrier constructions of Sections 14.1, 14.2 and 14.3 can be adapted to provide
boundary modulus of continuity estimates for $C^0(\overline{\Omega}) \cap C^2(\Omega)$ solutions of
equation (14.1). In particular, we observe that if in the hypotheses of any of the
gradient estimates of these sections, we only assume that $u \in C^0(\overline{\Omega}) \cap C^2(\Omega)$ then
we obtain, in place of an estimate for $\sup_{\partial\Omega} |Du|$, a bound for the quantity
\[
\sup_{\substack{x \in \Omega \\ y \in \partial\Omega}} \frac{|u(x) - u(y)|}{|x - y|}.
\]

Furthermore a boundary modulus of continuity estimate still results when we only
assume that the boundary values $\varphi \in C^0(\partial\Omega)$. To show this we fix a point $y$ on $\partial\Omega$
and for arbitrary $\varepsilon > 0$, choose $\delta > 0$ such that $|\varphi(x) - \varphi(y)| < \varepsilon$ whenever $|x - y| < \delta$.
We then define functions $\varphi^{\pm} \in C^2(\overline{\Omega})$ by

\begin{equation}
\varphi^{\pm}(x) = \varphi(y) \pm \left(\varepsilon + \frac{2 \sup |\varphi|}{\delta^2} |x - y|^2\right).
\tag{14.87}
\end{equation}

Clearly, on $\partial\Omega$ we have
\[
\varphi^- \leq \varphi \leq \varphi^+
\]

Hence if the operator $\Qop$ and the functions $\varphi^{\pm}$ satisfy any of the conditions of the
estimates derived in Sections 14.1, 14.2 and 14.3, we obtain an estimate

\begin{equation}
\varphi^- + w^- \leq u \leq \varphi^+ + w^+
\tag{14.88}
\end{equation}

in $\mathcal N \cap \Omega$ where $w^{\pm}$ are the relevant barrier functions and $\mathcal N$ is some neighborhood
of $y$. Since in all our previous barrier constructions we had $w^+ = -w^- = w$, it
follows from (14.88) that

\begin{equation}
|u(x) - \varphi(y)| \leq \varepsilon + w(x) + \frac{2 \sup |\varphi|}{\delta^2} |x - y|^2
\tag{14.89}
\end{equation}

in $\mathcal N \cap \Omega$. We can therefore assert the following continuity estimates.

\begin{theorem}[14.15]\label{gt14-boundary-gradient-14-15}\dcref{Theorem 14.15}\group{ch14-boundary-gradient}\source{gilbarg-trudinger:page-0365}
Let $u, \varphi \in C^0(\overline{\Omega}) \cap C^2(\Omega)$ satisfy $\Qop u = 0$ in $\Omega$ and $u = \varphi$ on $\partial\Omega$.
Suppose that the operator $\Qop$ and domain $\Omega$ satisfy the structural and geometric
conditions of either Theorem 14.1, Corollary 14.3, Corollary 14.5, Corollary 14.7 or
Theorem 14.9. Then the modulus of continuity of $u$ on $\partial\Omega$ can be estimated in terms of
the modulus of continuity of $\varphi$ on $\partial\Omega$, $\sup_{\partial\Omega} |\varphi|$, $\sup_{\Omega} |u|$, $\Omega$ and the coefficients of $\Qop$.
\end{theorem}

% PDF page 366
\source{gilbarg-trudinger:page-0366}

\setcounter{page}{354}

14.6.\ Appendix: Boundary Curvatures and the Distance Function

Let $\Omega$ be a domain in $\mathbb{R}^n$ having non-empty boundary $\partial\Omega$. The distance function $d$
is defined by

\begin{equation}
d(x)=\operatorname{dist}(x,\partial\Omega).
\tag{14.90}
\end{equation}

It is readily shown that $d$ is uniformly Lipschitz continuous. For let $x,\, y \in \mathbb{R}^n$ and
choose $z \in \partial\Omega$ such that $|y-z|=d(y)$. Then

\[
d(x)\le |x-z|\le |x-y|+d(y)
\]

so that by interchanging $x$ and $y$ we have

\begin{equation}
|d(x)-d(y)|\le |x-y|.
\tag{14.91}
\end{equation}

Now let $\partial\Omega \in C^2$. For $y \in \partial\Omega$, let $v(y)$, and $T(y)$ denote respectively the unit
inner normal to $\partial\Omega$ at $y$ and the tangent hyperplane to $\partial\Omega$ at $y$. The curvatures of
$\partial\Omega$ at a fixed point $y_0\in\partial\Omega$ are determined as follows. By a rotation of coordinates
we can assume that the $x_n$ coordinate axis lies in the direction $v(y_0)$. In some
neighborhood $\mathcal{N}=\mathcal{N}(y_0)$ of $y_0$, $\partial\Omega$ is then given by $x_n=\varphi(x')$ where $x'=
(x_1,\ldots,x_{n-1})$, $\varphi\in C^2(T(y_0)\cap\mathcal{N})$ and $D\varphi(y_0)=0$. The curvature of $\partial\Omega$ at $y_0$ is
then described by the orthogonal invariants of the Hessian matrix $[D^2\varphi]$ evaluated
at $y_0$. The eigenvalues of $[D^2\varphi(y_0)]$, $\kappa_1,\ldots,\kappa_{n-1}$, are called the principal curvatures
of $\partial\Omega$ at $y_0$ and the corresponding eigenvectors are called the principal directions
of $\partial\Omega$ at $y_0$. The mean curvature of $\partial\Omega$ at $y_0$ is given by

\begin{equation}
H(y_0)=\frac{1}{n-1}\sum \kappa_i=\frac{1}{n-1}\,\Delta\varphi(y_0).
\tag{14.92}
\end{equation}

By a further rotation of coordinates we can assume that the $x_1,\ldots,x_{n-1}$ axes lie
along principal directions corresponding to $\kappa_1,\ldots,\kappa_{n-1}$ at $y_0$. Let us call such a
coordinate system a principal coordinate system at $y_0$. The Hessian matrix
$[D^2\varphi(y_0)]$ with respect to the principal coordinate system at $y_0$ described above
is given by

\begin{equation}
[D^2\varphi(y_0)]=\operatorname{diag}[\kappa_1,\ldots,\kappa_{n-1}].
\tag{14.93}
\end{equation}

The unit inner normal vector $\bar v(y)=v(y)$ at the point $y=(y',\varphi(y'))\in\mathcal{N}\cap\partial\Omega$
is given by

\begin{equation}
v_i(y)=\frac{-D_i\varphi(y')}{\sqrt{1+|D\varphi(y')|^2}},\qquad i=1,\ldots,n-1,\qquad
v_n(y)=\frac{1}{\sqrt{1+|D\varphi(y')|^2}}.
\tag{14.94}
\end{equation}

Hence with respect to the principal coordinate system at $y_0$, we have

\begin{equation}
D_j\bar v_i(y_0)=-\kappa_i\,\delta_{ij},\qquad i,j=1,\ldots,n-1.
\tag{14.95}
\end{equation}


% PDF page 367
\source{gilbarg-trudinger:page-0367}

14.6. Appendix: Boundary Curvatures and the Distance Function \hfill 355

For $\mu>0$, let us set $\Gamma_{\mu}=\{x\in\curlyO\mid d(x)<\mu\}$. The following lemma relates the smoothness of the distance function $d$ in $\Gamma_{\mu}$ to that of the boundary $\partial\Omega$.

\begin{lemma}\label{gt14-boundary-gradient-14-16}\dcref{Lemma 14.16}\group{ch14-boundary-gradient}\source{gilbarg-trudinger:page-0367}
\textbf{Lemma 14.16.} \textit{Let $\Omega$ be bounded and $\partial\Omega\in C^k$ for $k\ge 2$. Then there exists a positive constant $\mu$ depending on $\Omega$ such that $d\in C^k(\Gamma_{\mu})$.}

\textit{Proof.} The conditions on $\Omega$ imply that $\partial\Omega$ satisfies a uniform interior sphere condition; that is, at each point $y_0\in\partial\Omega$ there exists a ball $B$ depending on $y_0$ such that $\overline{B}\cap(\R^n-\Omega)=y_0$ and the radii of the balls $B$ are bounded from below by a positive constant, which we take to be $\mu$. It is easy to show that $\mu^{-1}$ bounds the principal curvatures of $\partial\Omega$. Also, for each point $x\in\Gamma_{\mu}$, there will exist a unique point $y=y(x)\in\partial\Omega$ such that $|x-y|=d(x)$. The points $x$ and $y$ are related by

\begin{equation}
(14.96)\qquad x=y+\vnu(y)d.
\end{equation}

We show that this equation determines $\vnu$ and $d$ as $C^{k-1}$ functions of $x$. For a fixed point $x_0\in\Gamma_{\mu}$, let $y_0=y(x_0)$ and choose a principal coordinate system at $y_0$. We define a mapping $\mathbf{g}=(g^1,\dots,g^n)$ from $\mathcal U=(T(y_0)\cap \N(y_0))\times\R$ into $\R^n$ by

\[
\mathbf{g}(y',d)=y+\vnu(y)d,\qquad y=(y',\varphi(y')).
\]

Clearly $\mathbf{g}\in C^{k-1}(\mathcal U)$, and the Jacobian matrix of $\mathbf{g}$ at $(y'_0,d(x))$ is given by

\begin{equation}
(14.97)\qquad [D\mathbf{g}]=\diag[1-\kappa_1 d,\dots,1-\kappa_{n-1}d,1].
\end{equation}

Since the Jacobian of $\mathbf{g}$ at $(y'_0,d(x_0))$ is given by

\begin{equation}
(14.98)\qquad \det [D\mathbf{g}]=(1-\kappa_1 d(x_0))\cdots(1-\kappa_{n-1} d(x_0))>0\qquad \text{since } d(x_0)<\mu,
\end{equation}

it follows from the inverse mapping theorem that for some neighborhood $\M=\M(x_0)$, the mapping $y'$ is contained in $C^{k-1}(\M)$. From (14.96) we have $Dd(x)=\vnu(y(x))=\vnu(y'(x))\in C^{k-1}(\M)$ for $x\in\M$. Hence $d\in C^k(\M)$, and thus $d\in C^k(\Gamma_{\mu})$. $\square$
\end{lemma}

An expression for the Hessian matrix of $d$ at points close to $\partial\Omega$ is an immediate consequence of the proof of Lemma 14.16.

\begin{lemma}\label{gt14-boundary-gradient-14-17}\dcref{Lemma 14.17}\group{ch14-boundary-gradient}\source{gilbarg-trudinger:page-0367}
\textbf{Lemma 14.17.} \textit{Let $\Omega$ and $\mu$ satisfy the conditions of Lemma 14.16 and let $x_0\in\Gamma_{\mu}$, $y_0\in\partial\Omega$ be such that $|x_0-y_0|=d(x_0)$. Then, in terms of a principal coordinate system at $y_0$, we have}

\begin{equation}
(14.99)\qquad [D^2d(x_0)]=\diag\!\left[\frac{-\kappa_1}{1-\kappa_1 d},\dots,\frac{-\kappa_{n-1}}{1-\kappa_{n-1} d},0\right].
\end{equation}

\textit{Proof.} Since

\[
Dd(x_0)=\vnu(y_0)=(0,0,\dots,1)
\]

% PDF page 368
\source{gilbarg-trudinger:page-0368}

\pagestyle{empty}



we have $D_{in} d(x_0)=0$, $i=1,\ldots,n$. To obtain the other derivatives we write, for
$i,j=1,\ldots,n-1$,

\[
D_{ij}d(x_0)=D_j v_i\circ y(x_0)
=D_k \bar v_i(y_0)D_j y_k(x_0)
=
\begin{cases}
\dfrac{-\kappa_i}{1-\kappa_i d} & \text{if $i=j$}\\
0 & \text{if $i\neq j$}
\end{cases}
\]

by (14.95) and (14.97). \qed
\end{lemma}

Note that the result of Lemma 14.17 is equivalent to the geometrically evident
statement that the circles of principal curvature to $\partial\Omega$ at $y_0$ and to the parallel
surface through $x_0$ at $x_0$ are concentric.

\bigskip

\noindent\textbf{Mean Curvature}

\medskip

Let us derive now a formula for the mean curvature of a $C^2$ hypersurface $\Ssurf$ in
terms of its given representation. Let $y_0\in\Ssurf$ and suppose that in a neighborhood
$\Ncal$ of $y_0$, $\Ssurf$ is given by $\psi(x)=0$ where $\psi\in C^2(\Ncal)$ and $|D\psi|>0$ in $\Ncal$. The unit
normal to $\Ssurf$ at a point $y\in \Ssurf\cap\Ncal$ (directed towards positive $\psi$) is given by

\begin{equation}
\tag{14.100}
\mathbf{v}=\frac{D\psi}{|D\psi|}.
\end{equation}

\noindent Let $\kappa_1,\ldots,\kappa_{n-1}$ be the principal curvatures of $\Ssurf$ at $y_0$. Then with respect to a
corresponding principal coordinate system at y0, one can show that

\begin{equation}
\tag{14.101}
D_i\!\left(\frac{D_j\psi}{|D\psi|}\right)=-\kappa_i\,\delta_{ij}, \qquad i,j=1,\ldots,n-1,
\end{equation}

\[
D_i\!\left(\frac{D_n\psi}{|D\psi|}\right)=0,\qquad i=1,\ldots,n.
\]

Consequently the eigenvalues of the matrix $\bigl[D_j(D_i\psi/|D\psi|)\bigr]$ at $y_0$ with respect to
the original coordinates are $-\kappa_1,\ldots,-\kappa_{n-1},0$ and hence the mean curvature
of $\Ssurf$ at $y_0$ is given by

\begin{equation}
\tag{14.102}
H(y_0)=-\frac{1}{n-1}\left(D_i\left[\frac{D_i\psi}{|D\psi|}\right]\right)_{x=y_0}.
\end{equation}

In particular if $\Ssurf$ is the graph in $\mathbb{R}^{n+1}$ of a function of $n$ variables $u\in C^2(\Omega)$, that
is, $\Ssurf$ is defined by $x_{n+1}=u(x_1,\ldots,x_n)$, the mean curvature of $\Ssurf$ at $x_0\in\Omega$ is


% PDF page 369
\source{gilbarg-trudinger:page-0369}

\pagestyle{empty}


given by

\begin{equation}
\tag{14.103}
H(x_0)=\frac{1}{n}\left(D_i\left[\frac{D_i u}{\sqrt{1+|Du|^2}}\right]\right)_{x=x_0}.
\end{equation}

\noindent$\mathcal{G}$ is called a minimal surface if $H(x_0)=0$ for all $x_0\in\Omega$.

\noindent Note that we also obtain from (14.101) the following formula for the sum of the
squares of the principal curvatures $\kappa_1,\ldots,\kappa_n$ at $x_0$,

\begin{equation}
\tag{14.104}
\mathcal{G}^2=\sum_{i=1}^n \kappa_i^2(x_0)=\sum_{i,j=1}^n D_i v_j D_j v_i(x_0),
\end{equation}

\noindent where

\[
v_i=\frac{D_i u}{\sqrt{1+|Du|^2}}, \qquad i=1,\ldots,n.
\]

\noindent Finally, the Gauss curvature of $\mathcal{G}$ at $x_0$ is given by

\begin{equation}
\tag{14.105}
K(x_0)=\prod_{i=1}^n \kappa_i
=\det [D_i v_j]
=\frac{\det D^2u}{(1+|Du|^2)^{(n+2)/2}}.
\end{equation}

\bigskip

\noindent\textbf{Notes}

\medskip

Many of the basic features of our treatment of boundary gradient estimates are
already present in the early work of Bernstein for equations in two variables, [BE
1, 2, 3, 4, 5, 6]. Indeed, Bernstein employed auxiliary functions such as the function
$\psi$ in (14.4) for similar purposes and also considered the question of non-existence.
Bernstein's work in two variables was continued by Leray [LR], who also
considered the relationship between the solvability of the Dirichlet problem in a
domain $\Omega$ and the geometric nature of the boundary $\partial\Omega$. Finn showed that
convexity was a necessary and sufficient condition for the solvability of the
Dirichlet problem for the minimal surface equation in two variables [FN 2, 4].

The first striking results for equations in more than two variables were proved
by Jenkins and Serrin [JS 2], (see Theorem 14.14), and Bakelman [BA 5, 6]. A
general theory embracing many interesting examples was finally developed by
Serrin [SE 3] to which the results of Sections 14.3 and 14.4 are due. We have
adopted some minor simplifications from [TR 5] in the exposition of Section 14.3.


% PDF page 370
\source{gilbarg-trudinger:page-0370}

\noindent\textbf{Problems}

\medskip

\noindent\textbf{14.1.} Suppose that in the hypotheses of Theorem 14.6 we have $a_{\infty}^{ij}, b_{\infty}$ independent of $x$ and $\varphi \equiv 0$. Show that the estimate (14.52) continues to hold when the restriction on $\Lambda$ in (14.50) is removed.

\medskip

\noindent\textbf{14.2.} Suppose that only the coefficient $b$ of $Q$ is decomposed according to (14.43). Show that the result of Theorem 14.6 continues to hold provided we take $a_{0}^{ij}=0$ in (14.50) and replace the inequalities (14.51) by

\medskip

\begin{equation*}
(14.106)\qquad \kappa(y)\ge \pm b_{\infty}(y,\varphi(y),\pm v)
\end{equation*}

\noindent where $\kappa(y)$ is the minimum principal curvature of $\partial\Omega$ at $y$.

\medskip

\noindent\textbf{14.3.} Again, suppose that only the coefficient $b$ of $Q$ is decomposed according to (14.43). Show that the result of Theorem 14.9 continues to hold provided we take $a_{0}^{ij}=0$ in (14.59) and replace (14.60) by

\medskip

\begin{equation*}
(14.107)\qquad \kappa(y)> \pm b_{\infty}(y,\varphi(y),\pm v).
\end{equation*}

\medskip

\noindent\textbf{14.4.} Show that the results of Problems 14.2 and 14.3 can be sharpened by replacing the maximum eigenvalue $\Lambda$ by the trace $\mathcal{T}$ and compare these results with Theorem 14.5.

\medskip

\noindent\textbf{14.5.} Derive the assertion at the end of Section 14.4 (see [SE 3]).

% PDF page 371
\source{gilbarg-trudinger:page-0371}



\section*{Chapter 15}

\begin{center}
\LARGE Global and Interior Gradient Bounds
\end{center}

In this chapter we are mainly concerned with the derivation of apriori estimates
for the gradients of $C^2(\Omega)$ solutions of quasilinear elliptic equations of the form

\[
(15.1)\qquad Qu=a^{ij}(x,u,Du)D_{ij}u+b(x,u,Du)=0
\]

in terms of the gradients on the boundary $\partial\Omega$ and the magnitudes of the solutions.
The resulting estimates facilitate the establishment of Step III of the existence
procedure described in Section 11.3. On combination with the estimates of
Chapters 10, 13 and 14, they yield existence theorems for large classes of quasilinear
elliptic equations including both uniformly elliptic equations and equations of
form similar to the prescribed mean curvature equation (10.7). Since the methods
of this chapter involve the differentiation of equation (15.1), our hypotheses will
generally require structural conditions to be satisfied by the derivatives of the
coefficients $a^{ij}, b$. In Section 15.4 we shall see that these derivative conditions can be
relaxed somewhat for equations in divergence form, where different types of
arguments are appropriate.

We shall also consider in this chapter the derivation of apriori interior gradient
estimates. Such estimates lead to existence theorems for Dirichlet problems where
only continuous boundary values are assigned. Interior gradient bounds for
equations of mean curvature type will be treated in Chapter 16.

\section*{15.1. A Maximum Principle for the Gradient}

We commence with a gradient bound under relatively simple hypotheses which will
also serve as an illustration of the general technique to be applied in the following
section. Let us suppose that the principal coefficients of the operator $Q$ can be
written as follows:

\[
(15.2)\qquad a^{ij}(x,z,p)=a^{ij}_*(p)+\frac12\bigl(p_i c_j(x,z,p)+c_i(x,z,p)p_j\bigr)
\]

where $a^{ij}_*\in C^1(\mathbb{R}^n)$, $c_i\in C^1(\Omega\times\mathbb{R}\times\mathbb{R}^n)$, $i,j=1,\ldots,n$, and the matrix $[a^{ij}_*]$ is non-negative. The following are examples of such decompositions:
